%!TEX root = ../main.tex

\newpage
\thispagestyle{empty}
\vspace*{-2cm}
%%%
%%
%% Quatrième de couverture (last cover page of the memoirs)
%%
%%%

\hbox{\includegraphics[width=8.6cm]{ed_edmh-h.jpg}}


\bigskip
\noindent\fbox{\parbox{\textwidth}{
{\bf Titre : }  Sur l'apprentissage non supervisé en haute dimension 

\medskip
{\bf Mots Clefs : } clustering, agrégation, haute dimension, estimation de densité, mélanges.


\medskip
{\bf Résumé : }
\small
Deux sujet sont traités dans cette thèse: le clustering en haute dimension et l'estimation de densités de mélange.
L'estimation des paramètres d'une loi mélange est un problème difficile en haute dimension.
Trois méthodes sont présentées pour résoudre ce problème: la première est une estimation des matrices de covariances avec hypothèse de parcimonie, les deux autres visent à estimer le nombre de composantes du mélange.
La deuxième partie étudie l'estimateur du maximum de vraisemblance d'une densité sous l'hypothèse
qu'elle est bien approximée par un mélange de plusieurs densités données.
Nous réalisons une étude statistique des performances de l'estimateur par rapport à la perte de
Kullback-Leibler et établissons des bornes de risque sous la forme d'inégalités d'oracle exacte.
Nous introduisons la notion d’agrégation (presque)-D-parcimonieuse et des bornes inférieures sont établies.
Enfin, nous proposons un algorithme qui réalise l'agrégation en Kullback-Leibler de composantes d'un dictionnaire. Nous comparons sa performance avec différentes méthodes. Nous proposons ensuite une méthode pour construire le dictionnaire de densités et l’étudions de manière numérique.
\vspace{0.5cm}
}}

\bigskip
\noindent\fbox{\parbox{\textwidth}{
{\bf Title : }  On unsupervised learning in high dimension

\medskip
{\bf Keys words : }  clustering, aggregation, high dimension, density estimation, mixtures.

\medskip
{\bf Abstract : }
\small
Two subjects are treated in this thesis: high-dimensional clustering and estimation of mixture densities.
The estimation of the parameters of a mixture law is a difficult problem in high dimension.
Three methods are presented to solve this problem: the first is an estimation of covariance matrices with sparsity hypothesis, the other two are aimed at estimating the number of components of the mixture.
The second part studies the maximum likelihood estimator of a density under the assumption
that it is well approximated by a mixture of several given densities.
We perform a statistical study of the performance of the estimator with respect to the loss of
Kullback-Leibler and establish risk bounds in the form of exact oracle inequalities.
We introduce the concept of (nearly)-D-sparse aggregation and lower bounds are established.
Finally, we propose an algorithm that performs Kullback-Leibler aggregation of components of a dictionary. We compare its performance with different methods. We then propose a method to build the dictionary of densities and study it experimentally.
\vspace{0.5cm}
}}
\vspace{0.5cm}
\vfill
\hfill \includegraphics[width=1cm]{pictoParis-Saclay.jpg}
\end{document}
